\chapter*{ABSTRAK}
\addcontentsline{toc}{chapter}{ABSTRAK}

\begin{center}

    \textbf{RAG Terkondisi-Prediksi: Integrasi Prediksi Kadar Glukosa dan Rekomendasi Klinis Berbasis Dokumen untuk Pengelolaan Diabetes} \\
    \vspace{0.4cm}
    oleh \\
    \vspace{0.4cm}
    Daffari Adiyatma\\
    18222003\\
\end{center}


\begin{spacing}{1.0}
Sistem pendukung keputusan klinis untuk diabetes yang memakai prediksi glukosa maupun
\emph{Retrieval-Augmented Generation} (RAG) umumnya mengambil pengetahuan berdasarkan kondisi
pasien saat ini, sehingga rekomendasinya bersifat reaktif. Belum dijumpai penelitian yang
mengondisikan proses \emph{retrieval} pada kadar glukosa yang diprediksi. Tugas akhir ini
mengusulkan RAG terkondisi-prediksi, yaitu mekanisme yang menyusun kueri \emph{retrieval} dari
kadar glukosa hasil prediksi sehingga dokumen yang diambil dan rekomendasi yang dihasilkan bersifat
antisipatif, dalam alur yang keluarannya selalu diverifikasi dokter. Prediksi glukosa dilakukan
model \emph{Gradient Boosting} dengan fitur turunan fisiologis, sedangkan pengetahuan klinis
diambil dari 2.233 potongan pedoman melalui \emph{retrieval} leksikal. Evaluasi pada dataset
OhioT1DM memakai validasi silang lintas-pasien enam \emph{fold}. Pada horizon $+30$ menit, model
mencapai RMSE 20,72~mg/dL dengan 94,69\% titik pada zona A+B \emph{Clarke Error Grid}, sebanding
dengan \emph{Random Forest} dan LSTM; pada skenario pemantauan mandiri tanpa CGM (horizon sekitar
240 menit) galat naik menjadi 71,60~mg/dL, tetapi tetap di bawah \emph{baseline} persistence
(95,31~mg/dL). Pada kasus divergen, ketika kondisi mendatang berbeda dari kondisi saat ini, RAG
terkondisi-prediksi mengungguli RAG standar secara konsisten pada keenam \emph{fold} (\emph{Mean
Reciprocal Rank} 0,190 menjadi 0,284; $p=0{,}031$) dan pada 120 kasus nyata (0,186 menjadi 0,321;
$p<0{,}001$). Lengan \emph{oracle} mencapai MRR 0,778, sehingga sebagian besar selisih yang tersisa
berasal dari ketepatan prediksi, bukan dari mekanisme \emph{retrieval}. Keterbatasan yang sama
tampak pada sensitivitas hipoglikemia, yaitu 17,3\% dengan ambang atas hasil regresi dan 67,2\%
dengan pengklasifikasi kondisi langsung. Seluruh komputasi lokal berjalan tanpa GPU (median
0,339~detik, memori 342~MB) dan setiap rujukan dapat ditelusuri sampai halaman cetak pedomannya.
Evaluasi awal oleh dua belas tenaga kesehatan menilai rekomendasi sesuai pedoman dengan verifikasi
dokter, dengan skor \emph{System Usability Scale} 55,21 yang menunjukkan kegunaan sistem masih
perlu ditingkatkan.

\textbf{Kata kunci:} RAG terkondisi-prediksi, prediksi kadar glukosa, \emph{Gradient Boosting}, sistem pendukung keputusan klinis, diabetes.
\end{spacing}
